\chapter{Switches and Layer-1 Devices}\label{ch:nw:switches-and-layer-1-devices}

A store-and-forward switch cannot send a 1\,518-byte frame on before it has received all of it: at
\qty{10}{\giga\bit\per\second} that is \qty{1.2}{\micro\second} of waiting, whatever the switch does afterwards. A \omterm{def:nw:switches-and-layer-1-devices:l1}{layer-1 switch}
sold to trading firms forwards between its ports in \qty{4}{\nano\second}, with less than 100 picoseconds of jitter, and in the
words of its data sheet it does ``not buffer or queue data'': it never looks at the frame at all. Between the two sit the
cut-through switches, which look only at the first bytes. Choosing among them, port by port, is most of what designing the
network inside a trading cage means.

This chapter follows a frame through each kind of device, then through the passive devices that let a firm watch its own
traffic, and ends with three designs of the same small cage: the latency of each path, the devices each path crosses, and the
single devices whose failure would cut a server off.

\section{Store-and-forward against cut-through}

\begin{definition}[Store-and-forward and cut-through switching]\label{def:nw:switches-and-layer-1-devices:modes}
A switch using \emph{store-and-forward switching}\index{store-and-forward switching} receives a frame completely, checks its
frame check sequence, and only then starts to send it on the output port. A switch using \emph{cut-through
switching}\index{cut-through switching} starts to send a frame as soon as it has read the header fields it needs to choose the
output port, while the rest of the frame is still arriving.
\end{definition}

\begin{proposition}[What a switch adds]\label{prop:nw:switches-and-layer-1-devices:add}
Let a frame of $\ell_{\mathrm f}$ bytes arrive at $R_{\mathrm L}$ on a port whose output is idle and of the same rate, and let
$\tau$ be the switch's internal decision and transfer time. A store-and-forward switch adds
$8\ell_{\mathrm f}/R_{\mathrm L} + \tau$ to the frame's latency; a cut-through switch that reads the first $h$ bytes adds
$8h/R_{\mathrm L} + \tau$, independent of the frame's length. If the output port is busy, or faster than the input (it would
run out of bits in the middle of the frame), a cut-through switch must buffer the frame and falls back to store-and-forward for
it.
\end{proposition}

\begin{proof}
The last bit leaves the sender at $8\ell_{\mathrm f}/R_{\mathrm L}$ and reaches the far end one serialisation after the first
bit left the last device. A store-and-forward switch emits its first bit $\tau$ after its last bit arrived, so the serialisation
is paid twice, once into the switch and once out; a cut-through switch emits its first bit $\tau$ after the $h$-th byte arrived,
and the two serialisations overlap. If the output is busy the bits must wait; if it is faster than the input it would send
faster than the bits arrive, and a frame cannot pause in the middle, so the frame must first be stored.
\end{proof}

\begin{omfigure}
\begin{tikzpicture}[font=\footnotesize,fr/.style={draw,minimum height=0.45cm,anchor=west,inner sep=0pt}]
\node[anchor=east] at (-0.2,2.1) {in (both)};
\node[fr,fill=omDef!20,minimum width=5.0cm] at (0,2.1) {frame arriving};
\node[anchor=east] at (-0.2,1.2) {out, cut-through};
\node[fr,fill=omMeth!25,minimum width=5.0cm] at (1.3,1.2) {frame leaving};
\node[anchor=east] at (-0.2,0.3) {out, store-and-forward};
\node[fr,fill=omThm!20,minimum width=5.0cm] at (5.8,0.3) {frame leaving};
\draw[densely dashed,gray] (0.5,2.45) -- (0.5,-0.3);
\draw[densely dashed,gray] (5.0,2.45) -- (5.0,-0.3);
\draw[<->] (0.5,-0.25) -- node[below,font=\scriptsize]{header read, then $\tau$} (1.3,-0.25);
\draw[<->] (5.0,-0.25) -- node[below,font=\scriptsize]{$\tau$} (5.8,-0.25);
\draw[->] (0,-0.9) -- (11.4,-0.9) node[right]{time};
\end{tikzpicture}

\omcaption{The same frame through a cut-through and a store-and-forward switch. Both wait for the switch's internal time
$\tau$; the cut-through switch starts after the header (left dashed line), the store-and-forward switch after the last bit
(right dashed line), so it adds one whole serialisation of the frame.}\label{fig:nw:switches-and-layer-1-devices:modes}
\end{omfigure}

A store-and-forward switch drops a frame whose check sequence is wrong before it goes further; a cut-through switch has already
sent most of it when it finds out, and can only mark the end of the frame as bad. That is the whole price of cut-through, and
in a colocation cage, where links are short and errors rare, it is a price firms pay.

\begin{example}[One order, three switches]\label{ex:nw:switches-and-layer-1-devices:order}
A 100-byte order frame at \qty{10}{\giga\bit\per\second}. A store-and-forward switch with a \qty{500}{\nano\second} internal time
(a model value, not a data sheet's) adds $80 + 500 = \qty{580}{\nano\second}$; a cut-through switch that reads 64 bytes adds
\qty{51.2}{\nano\second} plus its internal time; a layer-1 device adds only its own few nanoseconds. For a 1\,518-byte frame
the store-and-forward switch's toll rises to \qty{1\,714}{\nano\second}, the others' do not change
(\cref{fig:nw:switches-and-layer-1-devices:curve}).
\end{example}

\section{Layer-1 devices: fan-out, multiplexing and patching}

\begin{definition}[Layer-1 switch, fan-out, multiplexer]\label{def:nw:switches-and-layer-1-devices:l1}
A \emph{layer-1 switch}\index{layer-1 switch} connects ports at the level of the electrical or optical signal: it regenerates
the bits of one input on the outputs configured for it without decoding frames, so it neither reads addresses nor queues. A
\emph{layer-1 fan-out}\index{layer-1 fan-out} is a configuration that copies one input to several outputs at once. A
\emph{layer-1 multiplexer}\index{layer-1 multiplexer} merges the frames of several inputs onto one output; unlike a fan-out it must
recognise frame boundaries and hold a frame that arrives while another is being sent.
\end{definition}

The fan-out is the natural device for market data: the venue's feed arrives once on the cross-connect and must reach every
server that trades on it, all at the same time, with nothing to decide. The multiplexer is the natural device for orders: many
servers, one order-entry handoff, and no routing decision beyond ``all to that port''. A multiplexer cannot be as fast as a fan-out,
because two servers can send at once.

\begin{proposition}[Contention at a multiplexer]\label{prop:nw:switches-and-layer-1-devices:mux}
Two frames of serialisation time $s$ reach a multiplexer's output at times differing by $D$. If $|D| < s$ the later frame waits
$s - |D|$; otherwise neither waits. If each server's response to a common event arrives with an independent normal delay of
standard deviation $\sigma$, $D$ is normal with standard deviation $\sigma\sqrt2$ and the later of two orders waits with
probability $\P(|D| < s) = 2\Phi\bigl(s/(\sigma\sqrt2)\bigr) - 1$.
\end{proposition}

\begin{proof}
The output sends one frame at a time: the frame that arrives second starts when the first has finished, $s$ after the first
started, or at its own arrival if that is later. The difference of two independent normal variables of variance $\sigma^2$ is
normal with variance $2\sigma^2$.
\end{proof}

\begin{example}[Two servers, one handoff]\label{ex:nw:switches-and-layer-1-devices:two}
A 100-byte order occupies \qty{96}{\nano\second} of line time at \qty{10}{\giga\bit\per\second}. If two of the firm's servers
answer the same market event with delays of standard deviation \qty{50}{\nano\second}, the second order waits with probability
$2\Phi(96/70.7) - 1 = 0.83$. The firm's own servers are competing with each other at the multiplexer at the moment that matters
most, and the one whose response was faster by a few nanoseconds is the one that goes first.
\end{example}

The chapter's tutorial runs this race for two, four and eight servers (\cref{fig:nw:switches-and-layer-1-devices:mux}). With four
servers answering within \qty{50}{\nano\second} of one another, an order waits \qty{93}{\nano\second} on average and the last one
\qty{186}{\nano\second}, more than the multiplexer's own \qty{39}{\nano\second}; spread the answers over a microsecond and the average
falls to \qty{6}{\nano\second}. The tighter a firm's servers are, the more they queue behind one another.

\begin{dated}{2026-09}{Device latencies from data sheets}\label{dat:nw:switches-and-layer-1-devices:devices}
Arista's 7130 Connect Series \omterm{def:nw:switches-and-layer-1-devices:l1}{layer-1 switches} forward between ports in \qty{4}{\nano\second} (\qty{6}{\nano\second} for the
96-port model), with less than \qty{100}{\pico\second} of jitter and the same latency for one-to-many mirroring. Arista's MetaMux
multiplexes several inputs onto one output in \qty{39}{\nano\second} on average (\qty{42}{\nano\second} on another hardware
model); ``when multiple packets arrive at the multiplexer at the same time'' it queues them in input buffers, and without
contention its latency varies by $\pm\qty{7}{\nano\second}$. Cisco's Nexus 3548-X cut-through switch quotes latencies ``as low as''
\qty{250}{\nano\second}, \qty{200}{\nano\second} in its warp mode, and replicates one port to any number of others in its warp SPAN
mode ``as low as'' \qty{50}{\nano\second}.
\end{dated}

\begin{omfigure}
\begin{tikzpicture}
\begin{axis}[width=11.5cm,height=5.4cm,ymode=log,xmin=0,xmax=1600,ymin=2,ymax=4000,
  xlabel={frame size (bytes)},ylabel={latency added (ns, log scale)},
  ytick={3,10,30,100,300,1000,3000},yticklabels={3,10,30,100,300,1\,000,3\,000},
  xtick={64,256,512,1024,1518},xticklabels={64,256,512,1\,024,1\,518},
  tick label style={font=\small},label style={font=\small},grid=major,grid style={gray!25},
  legend columns=4,legend style={font=\footnotesize,at={(0.5,-0.3)},anchor=north,draw=none,fill=none,
  /tikz/every even column/.append style={column sep=8pt}},table/col sep=comma]
\addplot[omThm,very thick,mark=*,mark size=1.4pt] table[x=frame,y=store_forward]{figdata/networks/02-switches-and-layer-1-devices/forwarding.csv};
\addplot[omMeth,very thick,mark=square*,mark size=1.4pt] table[x=frame,y=cut_through]{figdata/networks/02-switches-and-layer-1-devices/forwarding.csv};
\addplot[omExo,very thick,mark=triangle*,mark size=1.8pt] table[x=frame,y=mux]{figdata/networks/02-switches-and-layer-1-devices/forwarding.csv};
\addplot[omDef,very thick,mark=diamond*,mark size=1.8pt] table[x=frame,y=l1]{figdata/networks/02-switches-and-layer-1-devices/forwarding.csv};
\legend{store-and-forward (model),cut-through,multiplexer,layer-1 switch}
\end{axis}
\end{tikzpicture}

\omcaption{Latency each device adds to a frame at \qty{10}{\giga\bit\per\second}, egress idle, by frame size. Only the
store-and-forward switch grows with the frame (its internal time of \qty{500}{\nano\second} is a model value); the others use the
data-sheet figures of \cref{dat:nw:switches-and-layer-1-devices:devices}. Data: \texttt{fig\_cage.py}.}\label{fig:nw:switches-and-layer-1-devices:curve}
\end{omfigure}

A \omterm{def:nw:switches-and-layer-1-devices:l1}{layer-1 switch} is also a patch panel under software control: any port can be connected to any other, or to several, by a
configuration change instead of a technician moving a cable. Firms use that to move a server from one venue's line to another,
to swap a failed device for a spare, or to copy a link to a capture device for an hour.

\section{Taps and mirror ports}

A firm that wants to know what it sent and received, and when, must copy its traffic without disturbing it.

\begin{definition}[Network tap, port mirroring]\label{def:nw:switches-and-layer-1-devices:tap}
A \emph{network tap}\index{network tap} is a device inserted in a link that copies its traffic to monitoring ports; a passive
optical tap does it by splitting the light of each fibre, sending a fixed share (its split ratio) to the monitor, needs no power
and adds only the length of its fibre. \emph{Port mirroring}\index{port mirroring} is a switch's copying of the traffic of some of
its ports to another port, in software or in its forwarding hardware.
\end{definition}

A tap is on the link itself: it cannot drop, reorder or delay the copy relative to the original, and it keeps working when the
switch does not. Its costs are the light it takes (a 50/50 split halves the power reaching the far end, which a long cross-connect
may not afford) and a monitoring port for each direction. A mirror port is free, since the switch is already there, but it is a
switch port like any other: mirroring both directions of a busy \qty{10}{\giga\bit\per\second} link sends up to
\qty{20}{\giga\bit\per\second} to a \qty{10}{\giga\bit\per\second} port, which drops half of it when both directions are full, exactly in
the bursts one wanted to see. Capture devices (chapter~5) therefore sit on taps, and on layer-1 mirror outputs, which copy at \omterm{def:nw:networking-for-trading:ser}{line
rate}.

\section{Designing the network inside one cage}

\begin{definition}[Out-of-band management network, direct-attach cable]\label{def:nw:switches-and-layer-1-devices:oob}
An \emph{out-of-band management network}\index{out-of-band management network} is a separate network, with its own switches and
links, over which a firm administers its devices and servers, so that configuration, monitoring and remote access never share
a port or a queue with trading traffic, and still work when the trading network does not. A \emph{direct-attach
cable}\index{direct-attach cable} is a copper cable assembly with the transceivers built into its ends, used for short links
within a rack.
\end{definition}

\begin{method}[Designing a trading cage's network]\label{met:nw:switches-and-layer-1-devices:design}
\begin{enumerate}
\item List the handoffs: each venue's market-data lines (A and B), its order-entry connections, and the firm's own links out.
\item Give market data a replication device (a \omterm{def:nw:switches-and-layer-1-devices:l1}{layer-1 fan-out}, or a cut-through switch where filtering by group matters), with
  line A and line B on different devices so that one failure never silences both.
\item Give orders a path that crosses as few devices as possible: a \omterm{def:nw:switches-and-layer-1-devices:l1}{layer-1 multiplexer} or a cut-through switch per handoff, and a
  second handoff on a second device.
\item Tap every handoff, both directions, and bring the copies to the capture devices; put the time reference
  (chapter~4) in the same cage.
\item Keep management out of band.
\item Check the design by removing every device and every cable in turn: each server must still receive at least one line and
  reach at least one order handoff.
\end{enumerate}
\end{method}

\begin{omfigure}
\begin{tikzpicture}[font=\footnotesize,b/.style={draw,rounded corners=2pt,minimum height=0.6cm,align=center,inner sep=2pt},>=Stealth]
\node[b,fill=gray!10] (A) at (0,2.4) {handoff A};
\node[b,fill=gray!10] (B) at (0,1.2) {handoff B};
\node[b,fill=gray!10] (O) at (0,-0.4) {order handoff 1};
\node[b,fill=gray!10] (P) at (0,-1.6) {order handoff 2};
\foreach \n/\y in {tA/2.4,tB/1.2,tO/-0.4,tP/-1.6} {\node[b,fill=omExo!15,minimum width=0.9cm] (\n) at (2.7,\y) {tap};}
\node[b,fill=omDef!20] (la) at (5.2,2.4) {layer-1\\ fan-out};
\node[b,fill=omDef!20] (lb) at (5.2,1.2) {layer-1\\ fan-out};
\node[b,fill=omMeth!20] (m1) at (5.2,-0.4) {multiplexer};
\node[b,fill=omMeth!20] (m2) at (5.2,-1.6) {multiplexer};
\node[b,minimum height=1.6cm,minimum width=1.6cm] (s1) at (9.4,1.7) {server 1};
\node[b,minimum height=1.6cm,minimum width=1.6cm] (s2) at (9.4,-0.9) {server 2};
\node[b,fill=gray!8] (cap) at (2.7,-3.0) {capture and time reference};
\foreach \a/\b in {A/tA,B/tB,O/tO,P/tP,tA/la,tB/lb,tO/m1,tP/m2} {\draw[thick] (\a) -- (\b);}
\foreach \a in {la,lb} {\draw[->,omDef,thick] (\a.east) -- (s1.west); \draw[->,omDef,thick] (\a.east) -- (s2.west);}
\foreach \a in {m1,m2} {\draw[<-,omMeth,thick] (\a.east) -- (s1.west); \draw[<-,omMeth,thick] (\a.east) -- (s2.west);}
\draw[gray,dashed] (tP) -- (cap);
\end{tikzpicture}

\omcaption{The layer-1 design of the chapter's cage. Each venue line goes through a tap to its own fan-out, which copies it to
both servers; each server's orders go through two multiplexers to two order-entry handoffs. Every handoff is tapped (one tap's
copy drawn); management is out of band and not drawn. Removing any single device or cable leaves every server with a line and an
order path.}\label{fig:nw:switches-and-layer-1-devices:cage}
\end{omfigure}

\Cref{tab:nw:switches-and-layer-1-devices:designs} compares three designs of the same cage, two servers and the same cables
(30~metres of cross-connect to each handoff, 3~metres inside the cage), computed with \texttt{firm.cagenet}. The commodity design
puts two store-and-forward switches in series; the cut-through design gives each line its own switch; the layer-1 design is
\cref{fig:nw:switches-and-layer-1-devices:cage}.

\begin{table}[ht]
\centering\small
\begin{tabular}{lrrrl}
\toprule
Design & Feed in (ns) & Order out (ns) & Sum (ns) & Single points \\
\midrule
Store-and-forward switches & 1\,439.8 & 1\,430.2 & 2\,870.0 & both switches, three cables \\
Cut-through switches & 560.0 & 556.8 & 1\,116.7 & none \\
Layer-1 devices & 262.8 & 294.6 & 557.3 & none \\
\bottomrule
\end{tabular}
\caption{Three cage designs: latency from the first bit at the venue handoff to the last bit at the server for a 104-byte feed
frame, and from the server to the order handoff for a 100-byte order, at \qty{10}{\giga\bit\per\second}, egress idle. Data:
\texttt{nw\_cage.design\_table()} on \texttt{firm.cagenet}.}\label{tab:nw:switches-and-layer-1-devices:designs}
\end{table}

The layer-1 design's feed path is 263 nanoseconds, of which 176 are the 36 metres of fibre and 83 the frame's own
serialisation: the devices add 4. Its order path is slower than its feed path by the multiplexer's 39 nanoseconds less the
fan-out's 4, and slower still when two servers answer the same event. The next nanoseconds are no longer in the devices but in
the cables and in the servers.

\section{Tutorial: three cages}\label{tut:nw:switches-and-layer-1-devices:tut}

\begin{tutorial}
\textbf{Goal.} Model a cage's network, compute its paths and single points of failure, and race orders through a multiplexer.
\textbf{End state:} \cref{tab:nw:switches-and-layer-1-devices:designs} and \cref{fig:nw:switches-and-layer-1-devices:curve,fig:nw:switches-and-layer-1-devices:mux}.
\begin{tutsteps}
\item \textbf{Devices.} \texttt{nw\_cage.DEVICES} holds the data-sheet latencies and the one model value;
  \texttt{forwarding\_curve()} evaluates \cref{prop:nw:switches-and-layer-1-devices:add} per frame size.
\item \textbf{Cages.} \texttt{design(name)} builds each cage from \texttt{firm.cagenet} devices and cables; \texttt{Cage.path} finds the
  route and \texttt{breakdown} splits its latency (\cref{lst:nw:switches-and-layer-1-devices:breakdown}).
\item \textbf{Failures.} \texttt{Cage.single\_points(feeds, orders, servers)} removes each element in turn;
  \texttt{design\_table()} writes the table.
\item \textbf{Contention.} \texttt{mux\_contention(k, jitter)} races $k$ orders through one multiplexer over 20\,000 events
  (\cref{lst:nw:switches-and-layer-1-devices:mux}); \texttt{python fig\_cage.py} writes the charts' data.
\end{tutsteps}
\textbf{What to change next.} Replace the cross-connects by 150~metres and see which term of the table dominates; give each server
its own multiplexer input from a second network card and measure how contention changes when the firm's four fastest strategies
share two multiplexers instead of one.
\end{tutorial}

\omcode{code/firm/cagenet/firm_cagenet.py}{91}{101}{A path's latency: one serialisation, the fibre, and what each device adds.}\label{lst:nw:switches-and-layer-1-devices:breakdown}

\omcode{code/networks/02-switches-and-layer-1-devices/python/nw_cage.py}{77}{88}{Orders from $k$ servers through one multiplexer:
each starts when it has arrived and the previous one has finished.}\label{lst:nw:switches-and-layer-1-devices:mux}

\begin{omfigure}
\begin{tikzpicture}
\begin{axis}[width=11.5cm,height=5.2cm,xmode=log,xmin=20,xmax=2000,ymin=0,ymax=320,
  xlabel={standard deviation of the servers' response delays (ns, log scale)},ylabel={mean wait of an order (ns)},
  xtick={25,50,100,200,400,800,1600},xticklabels={25,50,100,200,400,800,1\,600},
  tick label style={font=\small},label style={font=\small},grid=major,grid style={gray!25},
  legend columns=3,legend style={font=\footnotesize,at={(0.5,-0.3)},anchor=north,draw=none,fill=none,
  /tikz/every even column/.append style={column sep=8pt}},table/col sep=comma]
\addplot[omDef,very thick,mark=*,mark size=1.4pt] table[x=jitter_ns,y=k2]{figdata/networks/02-switches-and-layer-1-devices/mux.csv};
\addplot[omMeth,very thick,mark=square*,mark size=1.4pt] table[x=jitter_ns,y=k4]{figdata/networks/02-switches-and-layer-1-devices/mux.csv};
\addplot[omThm,very thick,mark=triangle*,mark size=1.8pt] table[x=jitter_ns,y=k8]{figdata/networks/02-switches-and-layer-1-devices/mux.csv};
\legend{2 servers,4 servers,8 servers}
\end{axis}
\end{tikzpicture}

\omcaption{Simulation: $k$ servers each send one 100-byte order in answer to the same event through one
\qty{10}{\giga\bit\per\second} multiplexer; mean queueing delay of an order against the spread of the servers' response times,
20\,000 events. Tightly clustered answers queue behind one another. Data: \texttt{fig\_cage.py}
(seeded).}\label{fig:nw:switches-and-layer-1-devices:mux}
\end{omfigure}

\section{Build: the cage network model}\label{bld:nw:switches-and-layer-1-devices:cagenet}

\begin{build}
\bfield{Purpose} The firm's cage as data: what is connected to what, with which device, over how much fibre, so that every path's
latency and every single point of failure is computed rather than remembered. Chapter~5 taps it and chapter~29 prices it.
\bfield{Interface} \texttt{firm\_cagenet}: \texttt{Device(name, kind, fabric\_ns)} with kinds \texttt{handoff}, \texttt{server},
\texttt{l1}, \texttt{mux}, \texttt{cut-through}, \texttt{store-and-forward}, \texttt{tap}; \texttt{Cable(a, b, length\_m, n\_g)};
\texttt{Cage(devices, cables)} with \texttt{path}, \texttt{latency\_ns}, \texttt{breakdown}, \texttt{single\_points(feeds, orders,
servers)} and \texttt{untapped(handoffs)}.
\bfield{Rules} Latency runs from the first bit at the source to the last bit at the destination; every device's toll comes from
\texttt{firm.netsim.forward\_ns}; taps add only fibre; paths never cross a server or a handoff; an unknown device kind or a cable to
an unknown device is an error.
\bfield{Acceptance tests} \texttt{code/firm/cagenet/tests/}: the toll of each device kind on a one-device chain; that a path does not
cross a server; single points of failure and untapped handoffs on a hand-built cage.
\bfield{Stretch} Load: route the feeds of chapter~1 through the cage with \texttt{firm.netsim.egress} at every output and report the
worst queue; a spares plan, listing for each single point the patch a \omterm{def:nw:switches-and-layer-1-devices:l1}{layer-1 switch} would make to route around it.
\end{build}

\begin{omsources}
\item Arista, \emph{7130 Connect Series Layer 1 Switch Data Sheet}; Arista, \emph{7130 MetaMux} product page.
\item Cisco, \emph{Nexus 3548-X, 3524-X, 3548-XL and 3524-XL Switches Data Sheet}.
\item Garland Technology, ``Network TAP Split Ratios and Loss Budget'' (2014).
\end{omsources}

\section{Exercises}

\begin{exercise}[$\star$]\label{exo:nw:switches-and-layer-1-devices:1}
What does a store-and-forward switch with a \qty{300}{\nano\second} internal time add to a 512-byte frame at
\qty{25}{\giga\bit\per\second}? And a cut-through switch reading 64 bytes, with the same internal time?
\end{exercise}

\begin{exercise}[$\star$]\label{exo:nw:switches-and-layer-1-devices:2}
A cut-through switch receives a frame at \qty{10}{\giga\bit\per\second} for an output at \qty{25}{\giga\bit\per\second}, then one at
25 for an output at 10. In which case can it cut through, and why?
\end{exercise}

\begin{exercise}[$\star$]\label{exo:nw:switches-and-layer-1-devices:3}
Both directions of a link are mirrored to one port. The link carries \qty{6}{\giga\bit\per\second} each way at the busiest moment;
the mirror port runs at \qty{10}{\giga\bit\per\second}. What share of the copy is lost at that moment?
\end{exercise}

\begin{exercise}[$\star\star$]\label{exo:nw:switches-and-layer-1-devices:4}
Two servers answer an event with independent normal delays of standard deviation \qty{100}{\nano\second}; their 100-byte orders go
through one \qty{10}{\giga\bit\per\second} multiplexer. What is the probability that the later order waits?
\end{exercise}

\begin{exercise}[$\star\star$]\label{exo:nw:switches-and-layer-1-devices:5}
In the layer-1 design of \cref{tab:nw:switches-and-layer-1-devices:designs}, split the feed path's 262.8~ns into serialisation,
fibre and devices. Which would you shorten first, and how?
\end{exercise}

\begin{exercise}[$\star\star$]\label{exo:nw:switches-and-layer-1-devices:6}
Why does a fan-out need no buffer while a multiplexer does? What does that imply for their latency and jitter?
\end{exercise}

\begin{exercise}[$\star\star\star$]\label{exo:nw:switches-and-layer-1-devices:7}
\emph{Coding.} Add a third order handoff and a third multiplexer to the layer-1 design, send server~1's orders through one and
server~2's through the others, and recompute the mean wait of \cref{fig:nw:switches-and-layer-1-devices:mux} for four servers at
\qty{50}{\nano\second} when the four split two and two across two multiplexers.
\end{exercise}

\begin{exercise}[$\star\star\star$]\label{exo:nw:switches-and-layer-1-devices:8}
\emph{Find the flaw.} ``Line A and line B come into the same \omterm{def:nw:switches-and-layer-1-devices:l1}{layer-1 switch}, which fans each out to our servers. It is a single box,
but \omterm{def:nw:switches-and-layer-1-devices:l1}{layer-1 switches} have no software to crash, so there is no single point of failure.''
\end{exercise}

\section{Problem: Three Designs for One Cage}

\begin{problem}[{Weekend problem --- nanoseconds, failures and money}]\label{pb:nw:switches-and-layer-1-devices:1}
A firm is choosing among the three designs of \cref{tab:nw:switches-and-layer-1-devices:designs} for a cage that trades on a signal
from the venue's feed and answers with an order to the same venue. Suppose that the cut-through design costs 60\,000 a year more
than the commodity design, and the layer-1 design 150\,000 a year more (currency units; the figures are the problem's, not market
prices).

\medskip\noindent\textbf{Part I --- The paths.}
\begin{enumerate}
\item What is each design's trigger path, from the first bit of the feed frame at the handoff to the last bit of the order at the
  order handoff, counting feed in and order out?
\item How much does the cut-through design save over the commodity one, and the layer-1 design over the cut-through one?
\item In the commodity design, how much of the feed path is the two switches' serialisation?
\item How long is the 36~metres of fibre of the layer-1 feed path, in nanoseconds?
\end{enumerate}

\medskip\noindent\textbf{Part II --- Contention.}
\begin{enumerate}[resume]
\item Four of the firm's strategies share the layer-1 design's first multiplexer and answer an event within
  \qty{50}{\nano\second} (standard deviation) of each other. What is the mean wait of an order (from the chapter's simulation)?
\item And of the last of the four?
\item Does the design still beat the cut-through design for the last order?
\item How would you spread the four strategies?
\end{enumerate}

\medskip\noindent\textbf{Part III --- Failures.}
\begin{enumerate}[resume]
\item List the commodity design's single points of failure.
\item Why does the cut-through design have none although each line has only one switch?
\item What must be true of the venue's order handoffs for the layer-1 design to have none?
\item What happens to latency when one of the layer-1 design's fan-outs fails?
\end{enumerate}

\medskip\noindent\textbf{Part IV --- The verdict.}
\begin{enumerate}[resume]
\item What is the cost per nanosecond per year of moving from commodity to cut-through?
\item And of moving from cut-through to layer-1?
\item State the \emph{named result}: the trigger path of each design and the cost of each nanosecond saved.
\item For which kind of strategy is the layer-1 step worth its price?
\item What does a layer-1 design give up that a switch provides?
\item Where do the next nanoseconds come from once the devices add only tens?
\item Which checks from the method would you run before signing off the design?
\item In one sentence: what is the fastest device in a cage?
\end{enumerate}
\end{problem}

\section{Interview questions}

\begin{interviewq}[$\star$ \iqroles{developer}]\label{iq:nw:switches-and-layer-1-devices:1}
What is the difference between store-and-forward and \omterm{def:nw:switches-and-layer-1-devices:modes}{cut-through switching}, and when does a cut-through switch store and forward?
\end{interviewq}

\begin{interviewq}[$\star$ \iqroles{developer}]\label{iq:nw:switches-and-layer-1-devices:2}
What is a \omterm{def:nw:switches-and-layer-1-devices:l1}{layer-1 switch}, and what can it not do that an Ethernet switch can?
\end{interviewq}

\begin{interviewq}[$\star\star$ \iqroles{developer}]\label{iq:nw:switches-and-layer-1-devices:3}
Why do firms put orders through a multiplexer rather than a fan-out, and what goes wrong when several servers send at once?
\end{interviewq}

\begin{interviewq}[$\star\star$ \iqroles{developer}]\label{iq:nw:switches-and-layer-1-devices:4}
Tap or mirror port for latency measurement? Defend your choice.
\end{interviewq}

\begin{interviewq}[$\star\star\star$ \iqroles{developer, trader}]\label{iq:nw:switches-and-layer-1-devices:5}
Design the network for a cage with two venue lines, one order handoff and four servers. Walk me through the failure cases.
\end{interviewq}

\begin{interviewq}[$\star\star$ \iqroles{developer}]\label{iq:nw:switches-and-layer-1-devices:6}
Your cut-through switch shows higher latency at the open than at noon, for the same frames. Why?
\end{interviewq}
